\CGZoneHeader{example-image-a}{ZONA BOULDER PRINCIPAL}
\CGTextSingleColumn{}{Texto descritivo sobre o setor...}

\CGSectorHeader{guide_cyan}{BOULDER PRINCIPAL}
\begin{CGSectorRoutes}
    \CGRoute{Mestre das Ilusões}{Xb}{20m}{Fixa}{}[10b 4+2]
    \CGRoute{Estherminadora}{IXb}{20m}{Fixa}{}[\CGstar3 ESTHERMINADORA 9b 5]
    \CGRoute{Perdas e Danos}{IXc}{20m}{Fixa}{}
    \CGRoute{Invasão de Privacidade}{VIIb}{20m}{Fixa}{}[\CGstar3 Inicia em um V0 até a primeira chapa.
Segue pela esquerda com um descanso em um platô após a terceira chapa. 
Crux após costurar a costura fixa antes da fenda. 
Decida somente por rapel.]
    \CGRoute{Vovó Lurdinha}{VIIb}{15m}{Fixa}{}["Via esportiva com agarras marcadas"
Saida difícil, e segue entre as vias DarkSide e Invasão de privacidade, sem usar as agarras dessas vias. Via de Ary Pacheco (2025)]
    \CGRoute{Dark side}{VIIa}{20m}{Fixa}{}[\CGstar3 ]
    \CGRoute{Intrometida}{VIIa}{20m}{Fixa}{}
    \CGRoute{Grampeleta}{VIsup}{20m}{Fixa}{}[\CGstar3 Via classica de caiada. Segue uma curvatura (dobtas ptigmáticas) tipica da Gneiss. Informação do Guia: "Aqui você vai encontrar algumas das primeiras vias equipadas em Serra
Caiada, a via Grampeleta (via esportiva que possui grampos e chapeletas),
o 6Sup mais escalado e almejado de
todo o pico".]
    \CGRoute{Fungos}{VIsup}{20m}{TopRope}{}[Compartilha o mesmo Top/Reunião da via Grampeleta.]
    \CGRoute{Via do Zé}{VIsup}{15m}{Fixa}{}[\CGstar2 Saida delicada. Segue reto cruzando a via Grampeleta.]
    \CGRoute{Tiranovaca}{Vsup}{15m}{Fixa}{}[\CGstar1 Parada em P simples. Alternativamente, pode terminar seguindo a direita na reunião da Via do Zé. A via é 3+1 e posteriormente foi adocionada uma quarta chapa, facilitando o termino alternativo na via da direita.]
    \CGRoute{Lado A/B}{IV}{10m}{Fixa}{}[Originalmente via de ascensão em Solo. A via mais facil do bloco e garante acesso ao topo do bloco. O nome se refere ao lado direito com graduação de quarto grau, enquanto que a graduação pelo lado esquerda com crux de sexto/setimo para passar a barriga.]
    \CGRoute{Insônia}{VIIa}{10m}{TopRope}{}[Mesmo top da Lado A/B.]
    \CGRoute{Bolha D'água}{VI}{10m}{Fixa}{}[\CGstar2 ]
    \CGRoute{Ursinhos Gummy}{Vsup}{10m}{TopRope}{}
    \CGRoute{VáCliff}{AID:A1}{15m}{Fixa}{}
    \CGRoute{Tartaruga Ninja}{VIsup}{15m}{Fixa}{}[Via originalmente Top Rope, Posteriormente equipada.]
    \CGRoute{Alternativa de Mulambo}{IV}{15m}{Fixa}{}[\CGstar1 Via originalmente Top Rope, Posteriormente equipada.]
    \CGRoute{Mãe Solteira}{VIIa}{20m}{Fixa}{}
\end{CGSectorRoutes}

\CGSectorHeader{guide_cyan}{FORMIGA}
\begin{CGSectorRoutes}
    \CGRoute{Formiga no Crucifixo}{V}{15m}{Móvel}{}[Top Rope ou Móvel]
    \CGRoute{Formiga Nanica}{Vsup}{15m}{TopRope}{}[Top Rope]
    \CGRoute{Formiga Atômica}{VIIa}{15m}{TopRope}{}[Top Rope]
    \CGRoute{Formiga saúva}{VIsup}{15m}{TopRope}{}[Top Rope]
    \CGRoute{Formiga Rasgada}{V}{15m}{Móvel}{}[Móvel (sem Top Rope)]
    \CGRoute{Perdidos no Formigueiro}{IVsup}{15m}{TopRope}{}[Top Rope]
    \CGRoute{Formiga Louca}{IV}{15m}{Fixa}{}[Esportiva 3+2]
    \CGRoute{Camaleão}{III}{15m}{TopRope}{}[\CGstar1 Top Rope]
    \CGRoute{Formiga Cotó}{IV}{15m}{Fixa}{}[\CGstar1 Esportiva 2+2
Parada adicional na virada, após a fenda.]
    \CGRoute{Formiga Amarela}{VI}{15m}{TopRope}{}[\CGstar1 Top Rope]
    \CGRoute{Formiga Cega}{V}{15m}{TopRope}{}[Top Rope]
\end{CGSectorRoutes}

